\begin{frame}{자주 묻는 질문}
  \begin{itemize}
    \item \textbf{Q. 우리 팀 발표에 대한 리뷰가 너무 비판적이면?}
      비판적인 리뷰일수록 체크리스트의 목적에 가깝다 ---
      발표팀이 못 보고 지나친 문제를 찾아주는 것이지 인신공격
      이 아니다. 다만 리뷰어도 ``왜 그 지적이 타당한지''를
      \textbf{이 학기 개념으로 뒷받침}해야 한다.
    \item \textbf{Q. 시간이 부족하면 무엇을 우선?}
      \textbf{``방법론 적절성''}을 가장 먼저 --- 문제 정의가
      아무리 명확해도 데이터 특성에 안 맞는 모델(예: 불균형에
      정확도)을 골랐다면 뒤의 \textbf{결과 해석 전체}가
      흔들리기 때문.
    \item \textbf{Q. 발표팀이 ``좋은 지적입니다''만 하면?}
      그건 발표팀이 \textbf{아직 답하지 못했다는} 증거. 좋은
      리뷰는 ``좋은 지적''으로 넘길 수 \textbf{없는} 질문 ---
      답하려면 \textbf{새로운 계산·실험·그래프}를 꺼내야 하는
      질문. ``Precision·Recall을 보여주세요''는 그 자리에서
      혼동행렬을 다시 써야 답 $\to$ \textbf{반박 가능}.
      ``모델을 더 생각해보세요''는 어떤 답도 가능 $\to$ 반박
      불가능. ``좋은 지적''으로 끝나면 그 질문은 2점에 머문
      것 --- \textbf{구체적 산출물(어떤 그래프/계산)을 요구}하는
      쪽으로 재작성.
  \end{itemize}
\end{frame}
